% Checks for the section on learned interfaces (key lrn)
\subsection{Checks for \texorpdfstring{Section~\ref{sec:lrn}}{learned interfaces}}\label{chk:lrn}

\begin{numcheck}[Tail of the understanding error]\label{chk:lrn-understanding}
% src: lrn_understanding.py, part A
For \cref{thm:lrn-understanding-moment}, instances were drawn with a code $H$ whose singular values are evenly spaced in $[1,1.6]$, prior covariance $P=0.5I+0.1AA^\top$ for a standard Gaussian matrix $A$, channel covariance $R_c=0.05I+0.02\one\one^\top$ and code error scale $\sigma_H=0.7$. For each pair $(K,r_c)$, $4\times10^5$ draws of the understanding error $\psi$ were generated and the tail index of $\Vert\psi\Vert^2$ was estimated by the Hill estimator on the 2000 largest values; the moment threshold of the theorem corresponds to the index $d/2$ with $d=K-r_c+1$. The estimates were $0.469$ and $0.490$ against $0.5$ for $(K,r_c)=(2,2)$ and $(3,3)$, $1.037$ and $1.037$ against $1$ for $(3,2)$ and $(4,3)$, $1.448$ and $1.488$ against $1.5$ for $(4,2)$ and $(5,3)$, and $2.556$ against $2$ for $(5,2)$. The last estimate misses the predicted value by 28 percent; at this sample fraction the estimator is not reliable for the lighter tail, and the case $d=4$ rests on the proof alone. The running mean of $\Vert\psi\Vert^2$ over the first $10^3$, $10^4$, $10^5$ and $4\times10^5$ draws was $665$, $834$, $1.9\times10^5$, $3.9\times10^5$ at $(2,2)$ and $1.2\times10^4$, $6.6\times10^4$, $2.0\times10^7$, $5.6\times10^7$ at $(3,3)$, where $d=1$; it was $3.70$, $3.67$, $4.12$, $28.4$ at $(3,2)$ and $9.98$, $9.62$, $10.9$, $19.8$ at $(4,3)$, where $d=2$ and the second moment diverges logarithmically; and over the last three sample sizes it was $1.27$, $1.25$, $1.25$ at $(4,2)$, $19.7$, $19.0$, $18.9$ at $(5,3)$ and $0.446$, $0.448$, $0.450$ at $(5,2)$, where $d\ge3$. At small $\sigma_H$ the divergence is carried by rare draws in which the code error is comparable to the smallest singular value of $H$: at $K=3$, $r_c=2$ and $\sigma_H^2=0.02$, one draw of $4\times10^5$ had $\sigma_H\Vert W\Vert_F>0.9$ (\cref{chk:lrn-expansion}). A sample of $10^5$ draws at that scale usually contains no such draw and returns a stable average although the moment is infinite. For its instance with these dimensions \cref{chk:fw-efe} therefore reports no untruncated average and computes the moment truncated to an event $\sigma_H\Vert W\Vert_F\le\rho$ with $\rho$ below the smallest singular value of the code.
\end{numcheck}

\begin{numcheck}[Second-order expansion of the understanding error]\label{chk:lrn-expansion}
% src: lrn_understanding.py, part B; lrn_understanding2.py
For \cref{prop:lrn-understanding-expansion}, with the instance family of \cref{chk:lrn-understanding} and truncation radius $0.9$, below the smallest singular value of $H$, the truncated moment $\E[\psi\psi^\top;\sigma_H\Vert W\Vert_F\le0.9]-R_c$ was estimated with the channel noise integrated in closed form, from $4\times10^5$ draws ($3\times10^5$ in the last two cases). The ratio of the first term $\sigma_H^2\E\Vert\hat m_0\Vert^2$ to the diagonal entries of the estimate was $0.99$ and $0.95$ at $(K,r_c,\sigma_H^2)=(3,2,0.02)$, $1.00$ and $0.98$ at $(3,2,0.005)$, $1.08$ and $1.06$ at $(5,2,0.02)$, $1.08$ and $1.09$ at $(5,2,0.005)$, between $0.95$ and $0.97$ at $(4,4,0.005)$, and between $1.06$ and $1.15$ at $(6,3,0.01)$. The first term thus understates the excess in some instances and overstates it in others, and the overstatement at $(5,2)$ does not shrink with $\sigma_H$. Adding the three remaining terms of \cref{eq:lrn-understanding-expansion} reduced the largest entrywise gap in each case, in the same order, from $1.40\times10^{-3}$ to $1.08\times10^{-3}$, from $1.57\times10^{-4}$ to $7.8\times10^{-5}$, from $3.62\times10^{-3}$ to $1.75\times10^{-3}$, from $9.06\times10^{-4}$ to $1.28\times10^{-4}$, from $2.10\times10^{-3}$ to $1.55\times10^{-3}$, and from $4.41\times10^{-3}$ to $2.64\times10^{-3}$. The order of the remainder was tested separately with $10^5$ antithetic pairs $(W,-W)$, which cancel the odd orders. The mean of the second-order Taylor term agreed with the closed form within $1.8$ standard errors in all four instances, and the ratio of the largest entry of the mean remainder at $\sigma_H^2=0.004$ to that at $\sigma_H^2=0.001$ was $16.2$, $15.3$, $16.4$ and $15.7$ for $(K,r_c)=(3,2)$, $(5,2)$, $(4,4)$ and $(6,3)$, against $16$ for a remainder of order $\sigma_H^4$. The remainder is largest for the square code, about $50\sigma_H^4$ per diagonal entry at $K=r_c=4$.
\end{numcheck}

\begin{numcheck}[Least squares for a fixed code and tracking of a drifting one]\label{chk:lrn-tracking}
% src: lrn_tracking.py
For \cref{prop:lrn-code-risk}(a), the recursion \cref{eq:fw-rls} on the vectorized code and the recursion for $\Pi_t$ were run side by side for 12 rounds at $K=3$ and $r_c=2$: the covariances $P^H_t$ and $\Pi_t\otimes R_c$ differed by at most $1.0\times10^{-17}$ and the estimates by $1.4\times10^{-16}$. For part (c), with $M_2=\diag(1,0.5,0.25)$ and $4\times10^4$ replications, the mean of $\Vert\hat H_T-H\Vert_F^2$ at $T=5,6,8,20$ was $2.711$, $1.315$, $0.6633$, $0.1665$ against $2.664$, $1.332$, $0.6659$, $0.1665$ from \cref{eq:lrn-code-risk}. The standard errors at $T=8$ and $T=20$ were $0.0027$ and $0.0002$; for $T\le6$ the squared error has infinite variance and no standard error is quoted, so the averages at $T=5$ and $T=6$ are weaker evidence than those at $T=8$ and $T=20$. For \cref{thm:lrn-tracking}, one row with $K=3$, $\rho=0.05$, $M_2=\diag(1,0.5,0.25)$ and $P^-_0=0.1I$ was tracked by the exact recursion for $6\times10^4$ rounds, and averages were taken over the last four fifths. With $p_{\mathrm a}=1$ and $q=10^{-2}$, $10^{-3}$, $10^{-4}$, the first expression of \cref{eq:lrn-floor} was $0.1149$, $0.0328$, $0.01002$; the average of $\tr P^-_t$ was $0.1611$, $0.0373$, $0.01046$ and the realized squared error $0.1618$, $0.0378$, $0.01067$; the stationary error of the constant-gain rule used in the proof was $0.655$, $0.0420$, $0.01073$; and the right side of \cref{eq:lrn-tracking-upper} was $1.666$, $0.0444$, $0.0109$. With $p_{\mathrm a}=0.3$ and $q=10^{-3}$, $10^{-4}$, the floor was $0.0585$ and $0.01817$, the realized error $0.0789$ and $0.02072$, and the upper bound $0.1247$ and $0.02176$. The source script did not evaluate the stationary error of the constant-gain rule for $p_{\mathrm a}=0.3$; an independent recomputation gave $0.1084$ and $0.0211$ (\cref{chk:lrn-independent}), between the floor and the upper bound in both cases. The ratio of the average of $\tr P^-_t$ to the floor was $1.40$, $1.14$, $1.04$ at $p_{\mathrm a}=1$ and $1.31$, $1.10$ at $p_{\mathrm a}=0.3$, decreasing toward one with $q$ as parts (a) and (b) together require. With the third component of the regressor set to zero and $q=10^{-3}$, the entry $P^-_{33}$ after 2000 rounds was $2.1000$, equal to $p_0+qt$.
\end{numcheck}

\begin{numcheck}[Gauge of ungrounded learning]\label{chk:lrn-gauge}
% src: lrn_transfer.py, part A
For \cref{thm:lrn-gauge}, with $K=3$, a random orthogonal eigenbasis and the stationary content covariance $\Sigma=\sum_aq_a(1-f_a^2)^{-1}\Pi_a$, the Jacobian with respect to the code of the autocovariances $HF^k\Sigma H^\top$, $k=0,\dots,4$, was computed by central differences at a random code. Its nullity at $r_c=3$ and at $r_c=2$ was $3$ and $3$ for $F=0.9I$, $1$ and $1$ for eigenvalues $(0.9,0.9,0.5)$, and $0$ and $0$ for eigenvalues $(0.9,0.6,0.3)$, equal in each case to \cref{eq:lrn-orbit-dim} and consistent with \cref{prop:lrn-gauge-local}. For the simple spectrum the eight sign changes $T=V\diag(\pm1)V^\top$ altered the autocovariances by at most $1.2\times10^{-14}$, a generic rotation altered them by $4.99$, and at $r_c=3$ exactly one of the eight codes $HT$ agreed with $H$ on a grounded direction in general position, as \cref{cor:lrn-gauge-grounding} states.
\end{numcheck}

\begin{numcheck}[Composition along a path and hub holonomy]\label{chk:lrn-path}
% src: lrn_transfer.py, part B
For \cref{thm:lrn-path}(b), with $K=3$, errors $\Delta_e$ with independent $\Ncal(0,\delta^2/K)$ entries, $\delta^2=0.05$, a unit vector $x$ and $10^5$ replications, the mean of $\Vert(\Upsilon_{\mathfrak p}-I)x\Vert^2$ for $L=1,\dots,6$ was $0.0500$, $0.1027$, $0.1577$, $0.2150$, $0.2753$, $0.3384$, against $(1.05)^L-1=0.0500$, $0.1025$, $0.1576$, $0.2155$, $0.2763$, $0.3401$ from \cref{eq:lrn-path}. The standard errors grow from $10^{-4}$ at $L=1$ to $1.1\times10^{-3}$ at $L=6$, and the largest deviation is $1.6$ standard errors; the six averages share their draws and are not independent. For part (c), three hub maps with spoke errors of entry scale $0.3$ gave a triangle holonomy within $1.1\times10^{-14}$ of the identity.
\end{numcheck}

\begin{numcheck}[Risk decomposition, calibration cost and absorption]\label{chk:lrn-decomposition}
% src: lrn_transfer.py, parts C and D
For \cref{thm:lrn-risk-decomposition}(a), with $K=3$, two message rows, $M=\diag(1,2,0.5)$, a fixed code error with entries of scale $0.15$, $\bar h=(0.3,-0.2)^\top$, $\Sigma_h=0.02I$, the gain $\hat{\Kg}=P\hat H^\top(\hat HP\hat H^\top+R)^{-1}$ and $4\times10^5$ draws, the risk was $1.3544$ with standard error $0.0023$, against $0.8239+0.2589+0.1381+0.1340=1.3549$ for the channel, calibration, objective and interaction terms of \cref{eq:lrn-risk-decomposition}. The interaction term is positive in this instance and nearly as large as the objective term. For \cref{cor:lrn-calibration-cost}, with $\Pi=s\Pi_0$ for a fixed $\Pi_0\succ0$ and $2\times10^5$ antithetic pairs, the mean calibration cost was $0.006470$ (standard error $1.5\times10^{-5}$) at $s=0.02$ and $0.001618$ ($4\times10^{-6}$) at $s=0.005$, against $0.006468$ and $0.001617$ from the leading terms of \cref{eq:lrn-calibration-cost}. For \cref{prop:lrn-absorption}, $2\times10^5$ grounded rounds with content mean $0.7\one$, content covariance $\diag(1,0.64,0.36)$, a deviation matrix $D$ with entries of scale $0.3$, $d=(0.5,-0.4)^\top$ and $\Sigma_\zeta=0.04I$ gave a least-squares fit with intercept within $0.0019$ of $H+D$ and $0.0010$ of $d$, and a fit without intercept within $1.3\times10^{-3}$ of $H+D+d\bar m^\top M_2^{-1}$ evaluated at the sample moments.
\end{numcheck}

\begin{numcheck}[Independent recomputation]\label{chk:lrn-independent}
% src: lrn_verify_1_understanding.py, lrn_verify_2_tracking.py, lrn_verify_3_expansion.py, lrn_verify_4_gauge.py, lrn_verify_5_risk.py, lrn_verify_6_path.py, lrn_verify_7_conditional.py
The statements of this section were recomputed with separately written scripts, on separately drawn instances except for the tracking run, which uses the parameters of \cref{chk:lrn-tracking} with independently drawn regressors. For \cref{thm:lrn-understanding-moment}, at $\sigma_H=0.7$ with $(m_i,u,\nu)$ fixed and $4\times10^5$ draws of $W$, the slope of the logarithm of the empirical survival function of $\Vert\psi\Vert$ against $\log t$ was $-0.96$, $-2.39$, $-3.11$ at $(K,r_c)=(2,2)$, $(3,2)$, $(4,2)$ and $-1.01$, $-1.97$, $-3.16$ at $(3,3)$, $(4,3)$, $(5,3)$, against $-d$; the fit at $(3,2)$ is $20$ percent steeper than predicted. The empirical tail was at most $0.21$ times the bound \cref{eq:lrn-understanding-tail}. The Monte Carlo second moment was $2.65$, $1.72$, $1.95$ against $345$, $302$, $2069$ from \cref{eq:lrn-understanding-bound} at $(K,r_c,\sigma_H)=(4,2,0.7)$, $(6,2,0.7)$, $(7,3,0.3)$, so that bound holds with a wide margin and is far from tight. For \cref{prop:lrn-understanding-expansion}, with the channel noise integrated in closed form and antithetic pairs, the coefficient of $\sigma_H^2$ extrapolated from $\sigma_H=0.01$, $0.02$, $0.04$ agreed with the bracket of \cref{eq:lrn-understanding-expansion} within relative error $7\times10^{-3}$ at $(K,r_c)=(2,2)$, $(3,2)$, $(4,2)$, $(6,2)$, $(5,3)$, and the ratio of successive differences was between $3.97$ and $4.21$, against $4$ for a remainder of order $\sigma_H^4$. Over 2000 random instances at $K=3$ and $r_c=2$, the ratio of the first term of the bracket to its diagonal entries, computed from the closed form, had minimum $0.605$, fifth percentile $0.955$, median $0.999$, ninety-fifth percentile $1.047$ and maximum $25.1$; the extremes occur where a diagonal entry of the bracket is small. For \cref{rem:lrn-correction-sign}, the trace of the three noise terms was negative in all 1657 random instances with $K-r_c>\kappa(P)$, and for $P=pI$ it equaled $-2\sigma_H^2p(K-r_c-1)\tr(R_cS_0^{-1})$ to $1.1\times10^{-14}$.

For \cref{prop:lrn-code-risk}(c), with $K=3$, $M_2=\diag(1,0.5,0.25)$ and $\tr R_c=0.16$, the mean of $\Vert\hat H_T-H\Vert_F^2$ at $T=5,6,8,20$ was $1.147$, $0.558$, $0.279$, $0.0700$ against $1.120$, $0.560$, $0.280$, $0.0700$, and at $T=4$ the two halves of the sample gave $4.69$ and $6.61$, the instability expected of an infinite mean. For \cref{thm:lrn-tracking}, the first expression of \cref{eq:lrn-floor}, the stationary error of the constant-gain rule and the right side of \cref{eq:lrn-tracking-upper} agreed with \cref{chk:lrn-tracking} to the digits quoted there; the average of $\tr P^-_t$ was $0.1615$, $0.0374$, $0.01050$ at $p_{\mathrm a}=1$ and $0.0770$, $0.0200$ at $p_{\mathrm a}=0.3$; and the stationary error of the constant-gain rule at $p_{\mathrm a}=0.3$ was $0.1084$ and $0.0211$ from the fixed point and $0.1082$ and $0.0211$ in simulation. With an attention rule that depends on the receiver's history, 4000 replicas and 150 rounds, the smallest eigenvalue of $\E P^-_t-\bar P_t$ over all rounds was $5.6\times10^{-15}$, nonnegative as part (a) requires.

For \cref{thm:lrn-gauge} and \cref{prop:lrn-gauge-local}, the nullity of the differential of the autocovariances, computed from its closed form at $K=3$, was $3$, $3$, $2$ at $r_c=3,2,1$ for $F=0.9I$, was $1$ for eigenvalues $(0.9,0.9,0.5)$ and $0$ for $(0.9,0.6,0.3)$, equal to \cref{eq:lrn-orbit-dim}; at a rank-deficient block ($K=4$, multiplicities $(3,1)$, $r_c=2$, $\rank(HV_1)=1$) it was $4$, the orbit dimension $2$ plus the excess $2$ of \cref{prop:lrn-gauge-local}. At simple spectrum and $K=r_c=3$, a Gauss--Newton search from 60 random starts for codes with the autocovariances of a given code converged every time to one of the eight codes of the orbit and reached all eight. For \cref{cor:lrn-gauge-grounding}, at $K=r_c=4$ with multiplicities $(3,1)$, grounding of rank $0$, $1$, $2$, $3$ in general position left nullity $3$, $1$, $0$, $0$, the dimension of $\mathrm O(3-m_1)$, and at rank $2$ a reflection gave a second code that differs from $H$ by $3.06$ in its largest entry, agrees with $H$ on $\operatorname{col}(M_2)$ to $1.4\times10^{-15}$ and has the same autocovariances to $9.8\times10^{-15}$. With isotropic dynamics, $K=r_c=3$, $F=0.7I$ and $\rank M_2=2$, the grounded pairs alone leave a family of dimension $3$; with the stream the nullity was $0$, and the reflection across $\operatorname{col}(M_2)$ gave the one other code, at distance $2.36$ in the largest entry and equal to $H$ on $\operatorname{col}(M_2)$ to $2.3\times10^{-16}$. The stream therefore reduces the unidentified set under isotropic dynamics although it does not change the condition for point identification.

For \cref{thm:lrn-risk-decomposition}, with $K=3$, two message rows, a random code error with independent entries of standard deviation $0.35$, a regularized gain with $\hat R\ne R$ and $4\times10^5$ draws, the risk computed from the definitions was $3.5695$ with standard error $0.0068$, against $2.5326+0.7397+0.5243-0.2288=3.5678$ for the four terms of \cref{eq:lrn-risk-decomposition} computed from a separate sample; the interaction term is negative in this instance. The six weighted inner products of the four error vectors were within $1.3$ standard errors of zero. Two code errors with the same value of $\bar h-\Delta m_i$ gave offsets that differ by $0.177$ in norm, so the offset is not a function of $\bar h-\Delta m_i$ alone. Over $2\times10^4$ code models with scales spread over five orders of magnitude, the largest norm of the regularized gain was $2.2217$ against the bound $2.2384$ of part (d). For \cref{cor:lrn-calibration-cost}, the mean calibration cost was $2.819\times10^{-2}$, $6.972\times10^{-3}$, $1.738\times10^{-3}$ at $\Pi=s\Pi_1$ with $s=0.02$, $0.005$, $0.00125$, against $2.783\times10^{-2}$, $6.956\times10^{-3}$, $1.739\times10^{-3}$ from the leading terms. For \cref{eq:lrn-conditional-cost}, with the code model fixed and the true code random, the risk was $3.5753$ with standard error $0.0078$ against $3.5030+0.0680=3.5710$. For \cref{prop:lrn-absorption}, least squares with an intercept on $4\times10^5$ rounds was within $0.0027$ of $H+D$ and $0.0041$ of $d$, the fit without intercept was within $0.0018$ of $H+D+d\bar m^\top M_2^{-1}$, and the risk of the reading at the learned code was $1.8942$ with standard error $0.0039$ against $\tr(MP^+_D)=1.8944$, compared with $3.7804$ at the declared code.

For \cref{thm:lrn-path}, the conjugation identity of part (a) held to $3.8\times10^{-14}$ over 200 random walks; the mean of $\Vert(\Upsilon_{\mathfrak p}-I)x\Vert^2$ at $\delta^2=0.05$ and $L=1,\dots,6$ was $0.0499$, $0.1022$, $0.1572$, $0.2152$, $0.2759$, $0.3393$, against $0.0500$, $0.1025$, $0.1576$, $0.2155$, $0.2763$, $0.3401$ from \cref{eq:lrn-path}, with standard errors between $0.0001$ and $0.0008$ (the six averages come from nested walks and are not independent); hub holonomies over random closed walks were within $6.8\times10^{-15}$ of the identity with a consensus space of dimension $K=3$; the ratio of the two sides of \cref{eq:lrn-zero-shot} was at most $0.967$ over $2\times10^4$ random spoke pairs and equal to $1$ for scalar spokes $(1+\epsilon_i)I$ and $(1-\epsilon_k)I$, so the bound is attained; and independent pairwise translations with errors of scale $10^{-3}$ had a trivial consensus space in all 200 draws on a directed triangle, on the complete digraph with three agents and on a bidirected star with four agents. For \cref{rem:lrn-gaussian-spokes} the tail slope of the zero-shot mistranslation with Gaussian spokes was $-1.006$ against $-1$. For \cref{rem:lrn-two-maps} the mean-square mistranslation was $0.1336$ with standard error $0.0002$ against $0.1336$. For \cref{cor:lrn-hub-budget} the two thresholds agreed with a direct comparison of the two expressions for every $n$ from $2$ to $40$ on a grid of $c/T$ in $(0,3]$.
\end{numcheck}
